\documentclass[10pt,a4paper]{amsart}
\usepackage[margin=19mm]{geometry}
\usepackage{amsmath,amssymb,mathtools}
\usepackage[scr=rsfs,cal=euler]{mathalfa}
\usepackage{xfrac}
\usepackage[unicode,hidelinks,bookmarks=false]{hyperref}
\newcommand{\ie}{i.e.\ }
\newcommand{\cf}{cf.\ }
\newcommand{\resp}{respectively\ }
\newcommand{\Subsection}{\subsection}
\providecommand{\nobreakdash}{\nobreak\hbox{-}}
\setlength{\emergencystretch}{2em}
\multlinegap0pt
\usepackage{amssymb}
\usepackage{tikz}
\usepackage{cancel}
\usepackage[shortlabels]{enumitem}
\usepackage{stackengine}
\newtheorem{lem}{Lemma}
\newtheorem{prop}{Proposition}
\newcommand{\Kc}{\cali{K}}
\newcommand{\Om}{\Omega}
\newcommand{\al}{\alpha} 
\newcommand{\bR}{\mathbb{R}}
\newcommand{\cali}[1]{\mathscr{#1}} 
\newcommand{\tc}{{\tt c}}
\newcommand{\wed}{\wedge}
\newcommand{\wt}[1]{\widetilde{#1}}
\title{Alexander-Taylor's inequality for capacities in complex Sobolev spaces}
\author{Ngoc Cuong Nguyen, Do Duc Thai}
\date{Source: arXiv:2603.05838v1 (CC BY 4.0)}
\begin{document}
\maketitle

\noindent\emph{Source and licence.} Alexander-Taylor's inequality for capacities in complex Sobolev spaces. Version arXiv:2603.05838v1 (CC BY 4.0), distributed by arXiv under the Creative Commons Attribution 4.0 International licence, http://creativecommons.org/licenses/by/4.0/, as recorded in the arXiv licence metadata for this exact version and shown on the abstract page https://arxiv.org/abs/2603.05838v1. Full licence notice: \texttt{licenses/arxiv-2603.05838-CC-BY-4.0.txt}.\par\smallskip

\noindent\textit{Editorial scope.} Counted unit: the article's lem lem:extr-funciton-o (source0.tex lines 557--559). The article's complete original proof of it is delivered with it (source0.tex lines 561--580, one \texttt{proof} environment of 20 lines). No mathematics is written, repaired or re-derived: the statement and the proof are the article's own lines, in the article's own notation, and no retained line is substituted or re-flowed.  Retained uncounted as the source-local context the proof needs: lines 356--357; lines 360--362; lines 524--525.  Nothing was omitted from any retained span, so the package carries no omission record.  No retained line contains a citation, so the wrapper carries no bibliography and no bibliography entry.  No retained line contains a figure environment, an image-inclusion command or drawing code, so no image is delivered and the package declares no asset dependency.  Editorial formatting only, in addition: an AMS \texttt{amsart} preamble that restates the article's own declarations and macros used by the retained lines, namely the environments \texttt{lem}, \texttt{prop} on separate counters, together with the standard packages those lines need.  The retained environments print in this document as Lemma 1, Proposition 1 in order of appearance, whereas the article prints the same items with the numbering of its own full document.  The complete native source of the article as distributed in the arXiv e-print for this exact version is bundled unchanged as \texttt{sources/195825/source0.tex}.
\begin{lem} \label{lem:w-compactness}  Let $\{f_j \}_{j\geq 1} \subset W^*(X)$ be such that $\|f_j \|_* \leq A$ for every $j\geq 1$. There exists a subsequence $\{f_{j_k} \}_{j_k \geq 1}$ converging weakly in $W^{1,2}(X,\bR)$ to $f \in W^*(X)$  and $\|f\|_* \leq \liminf_{j\to \infty} \|f_j\|_*$.
\end{lem}

\begin{prop}  \label{prop:weak-conv-cp} Let $1<p<\infty$ and $E \subset \Om$ be a Borel set.  Let $\{f_j\} \subset L^p(\Om, \bR)$ and $\{g_j\} \subset L^p(\Om,\bR)$ be such that  $f_j \to f \in L^p(\Om,\bR)$ weakly and $g_j \to g\in L^p(\Om,\bR)$ weakly. 
If $|f_j(x)| \leq g_j(x)$ for a.e. $x\in E$, then $|f(x)| \leq g(x)$ a.e. $x\in E$. 
\end{prop}

\begin{prop} \label{prop:ae-qe}  Let $D\subset X$ be an open set and $f$ be a quasi-continuous function on $D$. If $f\geq 0$ a.e. on $D$, then $f\geq 0$  quasi-everywhere on  $D$.
\end{prop}

\begin{lem} \label{lem:extr-funciton-o} Assume $E\subset X$ is an open subset. There is a unique $\Phi_E \in W^*(X)$ such that 
$\tc (E) = \|\Phi_E\|_*^2$. Moreover, $-1\leq \Phi_E \leq 0$ a.e. and $\Phi_E=-1$ a.e. on $E$.
\end{lem}

\begin{proof} We first prove the existence. Let $\{u_j\}_{j\geq 1}$ be a sequence in $\Kc(E)$ such that $\tc(E) = \lim_{j\to \infty} \|u_j\|_*^2$. It follows from Lemma~\ref{lem:w-compactness} that $u_j$ converges weakly in $W^{1,2}(X)$ to  $u \in W^*(X)$ and $\|u\|_* \leq \liminf_j \|u_j\|_*$. Since $u_j \to u$ weakly in $L^2(X)$,  it follows from Proposition~\ref{prop:weak-conv-cp} that $u =-1$ a.e. in $E$ and $u\leq 0$ in $X$. Note that  the proof of that proposition used only the properties of Lebesgue points of functions in $L^2(X)$. Since $E$ is open,  $u\in \Kc(E)$ (this is the sole place that the openness of $E$ is used). By definition of  $\tc(E),$ we get $\tc(E) = \|u\|_*^2$. Here, we can choose $\Phi_E = \wt u$ which is quasi-continuous and $\Phi_E =-1$ q.e. on $E$.

Next,  we show the uniqueness. Suppose $u_1$ and $u_2$ are such two functions. Let $T_1$ and $T_2$ are the corresponding closed positive currents such that
$$
	\tc(E) = \|u_1\|_{L^2}^2 + \|T_1\| = \|u_2\|_{L^2}^2 + \|T_2\|.
$$  
Define $v = (u_1+u_2)/2$. Then, 
$
	dv \wed d^c v \leq (T_1 +T_2)/2 =:T.
$
So, $v\in \Kc(E)$ and 
$$\begin{aligned}
	\tc(E) &\leq \|v\|_{L^2}^2 + \|T\| \\&=\frac{1}{4} (\|u_1\|_{L^2}^2 + \| u_2\|_{L^2}^2) + \frac{1}{2} \|u_1u_2\|_{L^1} + \frac{1}{2} (\|T_1\|+ \|T_2\|). 
\end{aligned}$$
This implies that 
$$
	2\int_X u_1 u_2  \al^n \geq \int_X u_1^2 \;\al^n + \int_X u_2^2 \; \al^n.
$$
Thus, $u_1 = u_2$ a.e. If $u_1, u_2$ are quasi-continuous, then  $u_1=u_2$ q.e. by Proposition~\ref{prop:ae-qe}.
\end{proof}

\end{document}
